\documentclass[letterpaper,11pt]{article}

\usepackage[margin=1in]{geometry}
\setlength\parindent{0pt}
\setlength{\parskip}{0.7em}

\usepackage{titlesec}
\titlespacing*{\section}{0pt}{0.6\baselineskip}{0.2\baselineskip}
\titlespacing*{\subsection}{0pt}{0.4\baselineskip}{0.1\baselineskip}
\titlespacing*{\subsubsection}{0pt}{0.3\baselineskip}{0.1\baselineskip}

\usepackage{enumitem}
\setlist{nosep, topsep=2pt}

\usepackage{graphicx}
\graphicspath{{../results/interpretation/}}

\usepackage{subcaption}
\usepackage{booktabs}
\usepackage{amsmath}
\usepackage{amssymb}
\usepackage{bm}
\usepackage[hidelinks]{hyperref}

\title{Interpretation: SHAP and LIME Analysis --- Lab 3\\
       \large Connecting Language to fMRI, Stat 214, Spring 2026\vspace{-1em}}
\date{}

\begin{document}
\maketitle


\section{Interpretability Study}


Having established that the LoRA-fine-tuned BERT model with middle-layer pooling
(\texttt{finetuned\_lora\_mid4\_mean}) achieves the best voxelwise prediction
($\bar{r} = 0.0374$ for s2, $0.0512$ for s3), we ask \emph{what} each voxel is
responding to.  We apply two complementary attribution methods to the fitted
linear encoding model: closed-form linear SHAP identifies which BERT feature
dimensions drive predictions globally, and LIME isolates specific words within
individual stories that most influence particular voxels.  Both analyses are
conducted on the top-256 best-predicted voxels across two thematically distinct
test stories: \textit{On Approach to Pluto} (space exploration) and
\textit{Marry a Man Who Loves His Mother} (interpersonal relationships).
Predicted and observed BOLD response traces are shown for the four rank-1
voxels, with highlighted windows marking the local context used for LIME.


\subsection{Methods}


\subsubsection{SHAP for Linear Encoding Models}

The encoding model predicts voxel responses as $\hat{y}_i = \bm{X}_i \bm{w}$,
where $\bm{X}_i \in \mathbb{R}^{T \times 3072}$ is the delay-stacked embedding
matrix for story~$i$ and $\bm{w} \in \mathbb{R}^{3072}$ is the ridge-regression
weight vector for a given voxel.  Because $\bm{X}$ is z-scored before fitting,
the expected background value is zero and the linear SHAP attribution simplifies
to a closed form: the contribution of feature dimension~$j$ at time step~$t$ is
\begin{equation}
  \phi_j(x) = w_j (x_j - \mathbb{E}[x_j])
  \label{eq:shap}
\end{equation}
Global feature importance aggregates over the top-256 voxels:
\begin{equation}
  I_j = \sum_{v \in \mathcal{V}_{256}}
        \bigl|\phi_j^{(v)}\bigr|_{\text{mean}} \cdot \bigl|r_v\bigr|,
\end{equation}
where the outer weight $|r_v|$ up-weights voxels with stronger test correlations.
The 3072 feature dimensions correspond to four hemodynamic delay blocks
(delays 1--4\,TR $= 2$--8\,s), each of size 768 (the BERT hidden dimension).
We report importance both at the full feature level (3072 dims) and aggregated
across delay blocks to the 768 BERT dimension level.

Per-story word attribution maps each word to the TR in which it falls and
assigns the aggregate SHAP magnitude for that TR:
$\text{score}(w) = \sum_j |\phi_j(t_w)|$,
where $t_w$ is the TR index of word~$w$.  Because multiple words share the same
TR (at $\sim$3\,words/s with TR\,=\,2\,s, each TR contains up to six words),
words within the same TR receive identical scores.  This TR-resolution limitation
is addressed by LIME, which operates at the individual-word level.

\subsubsection{LIME with Word Masking}

LIME~(Ribeiro et al., 2016) approximates the local behavior of the encoding
model around a given input story by fitting a sparse linear surrogate on
perturbed inputs.  We represent each story as a binary vector over its unique
word types and generate 300 perturbations by randomly masking subsets of words.
Masking is implemented by zeroing all delay-block contributions of a word's
TR in the pre-built design matrix $\bm{X}$, preserving the original temporal
structure of surrounding context.  The ridge model then predicts the response
for each perturbed design matrix, and a weighted ridge regression on the binary
perturbation vectors fits the local linear surrogate.  The resulting
coefficients give each word a positive or negative contribution to the predicted
BOLD response of a target voxel.

The local window for LIME is defined as the 5-TR peak-response segment
(yellow shaded region in the response trace figures), identified as the
contiguous window with the highest mean absolute predicted BOLD signal during
that story.  Surrogate predictions average the model output across this
window, so LIME coefficients reflect how each word shifts the response
specifically during the most informative portion of the story.

We apply LIME to the three highest-ranked voxels (by test correlation $r$) for
each story, reporting the top positive contributors.


\subsection{SHAP Results}


\subsubsection{Dominant BERT Dimensions}

Figure~\ref{fig:shap_bertdim} shows the top-20 BERT dimensions by aggregate
SHAP importance for each subject.  The importance is broadly distributed across
all 768 dimensions with no single dominant axis, reflecting the high
dimensionality of the contextual embeddings.

\begin{figure}[H]
  \centering
  \begin{subfigure}[b]{0.48\textwidth}
    \includegraphics[width=\textwidth]{s2/finetuned_lora_mid4_mean/shap/s2_finetuned_lora_mid4_mean_shap_bert_dim_top20.png}
    \caption{Subject s2}
  \end{subfigure}
  \hfill
  \begin{subfigure}[b]{0.48\textwidth}
    \includegraphics[width=\textwidth]{s3/finetuned_lora_mid4_mean/shap/s3_finetuned_lora_mid4_mean_shap_bert_dim_top20.png}
    \caption{Subject s3}
  \end{subfigure}
  \caption{Top-20 BERT dimensions by aggregated SHAP importance for the
           \texttt{finetuned\_lora\_mid4\_mean} model.  Each bar is the mean
           absolute SHAP value across the top-256 voxels, weighted by test
           correlation.  The importance is distributed across many dimensions,
           with no single axis dominating.}
  \label{fig:shap_bertdim}
\end{figure}

Despite the distributed pattern, several dimensions recur across subjects.
BERT dimension~586 ranks 2nd for both s2 and s3, and dimension~539 appears in
the top-10 for both (Table~\ref{tab:top_bertdims}).  The importance values for
s2 are consistently higher (0.050--0.075) than for s3 (0.049--0.061), consistent
with the overall higher predictability of s3 spreading encoding signal across a
broader set of dimensions.

\begin{table}[ht]
  \centering
  \small
  \begin{tabular}{lrrr}
    \toprule
    & \multicolumn{2}{c}{Rank} & \\
    \cmidrule(lr){2-3}
    BERT dim & s2 & s3 & Shared \\
    \midrule
    586 & 2 & 2 & \checkmark \\
    539 & 8 & 5 & \checkmark \\
    160 & 1 & -- & \\
    19  & 2 & -- & \\
    545 & -- & 1 & \\
    536 & -- & 2 & \\
    \bottomrule
  \end{tabular}
  \caption{Selected BERT dimensions from the top-10 SHAP ranking.  Dimensions
           shared across subjects (586, 539) represent candidate universally
           brain-relevant semantic axes.}
  \label{tab:top_bertdims}
\end{table}



Examining importance at the full 3072-dimensional feature level reveals a
consistent pattern: for both subjects, the highest-ranked features overwhelmingly
belong to delay block~1 (1\,TR\,=\,2\,s lag).  Of the top-10 features for s2,
eight are from delay block~1; for s3, nine of ten.  The same BERT dimension
(e.g., dim~160 for s2) contributes at delay block~1 (rank~1) and again at
delay block~2 (rank~4), but with lower importance at the longer lag.

This dominance of the 2\,s lag is somewhat surprising given that the canonical
hemodynamic response function (HRF) peaks at approximately 5--6\,s.  It likely
reflects two factors: (1) ridge regression fits all four delays simultaneously
and can assign weight to early onsets of the neural response, and (2) naturalistic
narrative language exhibits strong temporal autocorrelation, so nearby lags carry
overlapping semantic information that inflates the apparent weight of short delays.



Figure~\ref{fig:shap_words} shows the top-20 words by SHAP score for each of the
two target stories for subject s2.  Because SHAP attribution operates at TR
resolution, words that co-occur within the same 2\,s window receive identical
scores; tied bars therefore reflect a shared high-salience time window rather
than individual word salience.

For the Pluto story, the dominant TRs contain age and temporal markers
(``year'', ``old'', ``five'', ``three'') clustered around a passage describing
the multi-year journey, along with motion verbs (``move'').  For the relationship
story, the dominant TR contains a dense cluster of physical-scene nouns
(``wooden'', ``chair'', ``propped'', ``foot'', ``against''), reflecting a
vivid spatial description that drives a concentrated burst of cortical activity.

\begin{figure}[H]
  \centering
  \begin{subfigure}[b]{0.48\textwidth}
    \includegraphics[width=\textwidth]{s2/finetuned_lora_mid4_mean/shap/s2_finetuned_lora_mid4_mean_onapproachtopluto_shap_words_top20.png}
    \caption{\textit{On Approach to Pluto}, s2}
  \end{subfigure}
  \hfill
  \begin{subfigure}[b]{0.48\textwidth}
    \includegraphics[width=\textwidth]{s2/finetuned_lora_mid4_mean/shap/s2_finetuned_lora_mid4_mean_marryamanwholoveshismother_shap_words_top20.png}
    \caption{\textit{Marry a Man\ldots}, s2}
  \end{subfigure}
  \caption{Top-20 words by SHAP score for subject s2.  Words sharing the same bar height belong to the same TR (2\,s window).  }
  \label{fig:shap_words}
\end{figure}


\subsection{LIME Results and Response Traces}


LIME resolves word-level attribution for the three top-ranked voxels per story.
Positive coefficients indicate words whose presence increases the voxel's
predicted BOLD response.  Response trace figures accompany the rank-1 voxels,
showing observed (blue) and predicted (orange) BOLD signals over the full story
with the yellow shaded region marking the 5-TR peak window used as the LIME
prediction target.

\subsubsection{Story: \textit{On Approach to Pluto} (Space Exploration)}

This story narrates the New Horizons spacecraft's multi-year journey to Pluto,
culminating in the July 2015 flyby.

\begin{figure}[H]
  \centering
  \begin{subfigure}[b]{0.48\textwidth}
    \includegraphics[width=\textwidth]{s2/finetuned_lora_mid4_mean/lime/onapproachtopluto/voxel63900_rank1_response.png}
    \caption{s2, rank-1 voxel (63900) --- response trace}
  \end{subfigure}
  \hfill
  \begin{subfigure}[b]{0.48\textwidth}
    \includegraphics[width=\textwidth]{s3/finetuned_lora_mid4_mean/lime/onapproachtopluto/voxel28817_rank1_response.png}
    \caption{s3, rank-1 voxel (28817) --- response trace}
  \end{subfigure}
  \caption{Observed (blue) and predicted (orange) BOLD response for the
           rank-1 voxel per subject during \textit{On Approach to Pluto}.
           The yellow shaded region marks the 5-TR peak window where LIME
           attributions are computed.}
  \label{fig:resp_pluto}
\end{figure}

\begin{figure}[H]
  \centering
  \begin{subfigure}[b]{0.48\textwidth}
    \includegraphics[width=\textwidth]{s2/finetuned_lora_mid4_mean/lime/onapproachtopluto/voxel63900_rank1_lime.png}
    \caption{s2, rank-1 voxel (63900)}
  \end{subfigure}
  \hfill
  \begin{subfigure}[b]{0.48\textwidth}
    \includegraphics[width=\textwidth]{s3/finetuned_lora_mid4_mean/lime/onapproachtopluto/voxel28817_rank1_lime.png}
    \caption{s3, rank-1 voxel (28817)}
  \end{subfigure}
  \caption{LIME word importances for \textit{On Approach to Pluto}, rank-1 voxel
           per subject.  Attributions correspond to the highlighted window in
           Figure~\ref{fig:resp_pluto}.}
  \label{fig:lime_pluto_rank1}
\end{figure}



The response traces (Figure~\ref{fig:resp_pluto}) show that the model tracks the
broad temporal envelope of the BOLD signal, with the predicted trace following
the shape of the observed response.  The peak window falls in the mid-story
portion where the narration shifts from scene-setting to active mission
description.  Prediction quality is visibly higher during the highlighted peak
than at quiet story segments, consistent with the voxels being selectively driven
by dense semantic content rather than baseline speech.

Across both subjects and both top voxels, the leading LIME words are temporal
markers and goal-directed action verbs tied to the mission narrative:
\texttt{july}, \texttt{fourth}, \texttt{decided}, \texttt{horizons},
\texttt{working}, \texttt{long}.  The word \texttt{i} (the narrator's
first-person voice) is prominent for s2, likely reflecting moments of personal
reflection that co-occur with high-salience mission content.  The word
\texttt{fourth} appears in both subjects' rank-2 voxels alongside \texttt{july},
confirming that the same passage drives activity in spatially neighbouring
voxels.

\subsubsection{Story: \textit{Marry a Man Who Loves His Mother} (Interpersonal Relationships)}

This story depicts close interpersonal dynamics, domestic settings, and social
negotiations between characters.

\begin{figure}[H]
  \centering
  \begin{subfigure}[b]{0.48\textwidth}
    \includegraphics[width=\textwidth]{s2/finetuned_lora_mid4_mean/lime/marryamanwholoveshismother/voxel32384_rank1_response.png}
    \caption{s2, rank-1 voxel (32384) --- response trace}
  \end{subfigure}
  \hfill
  \begin{subfigure}[b]{0.48\textwidth}
    \includegraphics[width=\textwidth]{s3/finetuned_lora_mid4_mean/lime/marryamanwholoveshismother/voxel25884_rank1_response.png}
    \caption{s3, rank-1 voxel (25884) --- response trace}
  \end{subfigure}
  \caption{Observed and predicted BOLD response for the rank-1 voxel per subject
           during \textit{Marry a Man Who Loves His Mother}.  The yellow region
           marks the 5-TR peak window used by LIME.}
  \label{fig:resp_marry}
\end{figure}

\begin{figure}[H]
  \centering
  \begin{subfigure}[b]{0.48\textwidth}
    \includegraphics[width=\textwidth]{s2/finetuned_lora_mid4_mean/lime/marryamanwholoveshismother/voxel32384_rank1_lime.png}
    \caption{s2, rank-1 voxel (32384)}
  \end{subfigure}
  \hfill
  \begin{subfigure}[b]{0.48\textwidth}
    \includegraphics[width=\textwidth]{s3/finetuned_lora_mid4_mean/lime/marryamanwholoveshismother/voxel25884_rank1_lime.png}
    \caption{s3, rank-1 voxel (25884)}
  \end{subfigure}
  \caption{LIME word importances for \textit{Marry a Man Who Loves His Mother},
           rank-1 voxel per subject.  Attributions correspond to the highlighted
           window in Figure~\ref{fig:resp_marry}.}
  \label{fig:lime_marry_rank1}
\end{figure}

\begin{figure}[H]
  \centering
  \begin{subfigure}[b]{0.48\textwidth}
    \includegraphics[width=\textwidth]{s2/finetuned_lora_mid4_mean/lime/marryamanwholoveshismother/voxel32271_rank2_lime.png}
    \caption{s2, rank-2 voxel (32271)}
  \end{subfigure}
  \hfill
  \begin{subfigure}[b]{0.48\textwidth}
    \includegraphics[width=\textwidth]{s3/finetuned_lora_mid4_mean/lime/marryamanwholoveshismother/voxel25885_rank2_lime.png}
    \caption{s3, rank-2 voxel (25885)}
  \end{subfigure}
  \caption{LIME word importances for \textit{Marry a Man Who Loves His Mother},
           rank-2 voxel per subject.}
  \label{fig:lime_marry_rank2}
\end{figure}

The response traces (Figure~\ref{fig:resp_marry}) place the highlighted window
in the latter portion of the story, where social negotiation between characters
reaches its narrative climax.  The model predictions closely follow the observed
signal in this region, again illustrating that prediction accuracy is
concentrated in semantically rich passages.

An inter-subject divergence emerges from the LIME attributions.  For s2,
high-importance words are relational verbs and adverbs describing co-habitation
and joint decisions: \texttt{miles}, \texttt{together}, \texttt{live},
\texttt{decided}, \texttt{asked}.  For s3, the leading terms are more concrete
physical-proximity and social-boundary words: \texttt{door}, \texttt{hand},
\texttt{came}, \texttt{let}, \texttt{soldier}, \texttt{wouldn't}.

Notably, the rank-1 and rank-2 voxels for each subject are immediate neighbours
in voxel index (s2: 32384 and 32271; s3: 25884 and 25885), yet their LIME
profiles diverge at a fine resolution while sharing the same social-proximity
theme.  This within-subject consistency is consistent with a smooth gradient of
semantic tuning in higher temporal cortex.


\subsection{Comparison of SHAP and LIME}

\subsubsection{Interpreted Results Layout}
SHAP and LIME answer related but distinct questions about the encoding model and
are best understood as complementary lenses rather than competing methods.
Table~\ref{tab:shap_vs_lime} summarises their key properties.

\begin{table}[ht]
  \centering
  \small
  \begin{tabular}{lll}
    \toprule
    & \textbf{SHAP} & \textbf{LIME} \\
    \midrule
    Scope        & Global across all test stories & Local to one story + one voxel \\
    Granularity  & BERT dim (768) + delay block   & Individual word token \\
    Attribution  & Closed-form: $\phi_j = X_{tj} w_j$ & Local linear surrogate (300 samples) \\
    Background   & Zero (z-scored features)       & Word masked to zero \\
    Time resolution & TR-level (2\,s windows)    & Word-level \\
    Output       & Feature importance across voxels & Word importance for peak window \\
    \bottomrule
  \end{tabular}
  \caption{Comparison of SHAP and LIME attribution methods as applied to the
           linear fMRI encoding model.}
  \label{tab:shap_vs_lime}
\end{table}

\subsubsection{Scope and Granularity}

SHAP operates globally: it aggregates over the entire test set and all 256
voxels to identify which BERT dimensions are systematically important across
stories and brain regions.  Its closed-form formula~(\ref{eq:shap}) requires no
sampling and is exact for linear models.  The cost is temporal resolution ---
because importance is averaged over all time steps, SHAP cannot resolve which
moments within a story drive the response.

LIME inverts this trade-off.  By constructing a local surrogate around the 5-TR
peak window of a single voxel's response, it achieves individual-word precision.
However, the 300-perturbation budget introduces sampling variance, the result is
specific to one story window and one voxel, and the local linear approximation
may not generalise to the broader response.

\subsubsection{Convergence on Key Content}

Despite their different scopes, SHAP and LIME converge on the same high-level
story content for both test stories.  For the Pluto story, SHAP identifies
TRs containing time-related vocabulary (``year'', ``old'', ``move'') as globally
important, and LIME independently assigns the highest word-level weights to
temporal mission markers (\texttt{july}, \texttt{fourth}, \texttt{working},
\texttt{long}) within the peak window.  For the relationship story, SHAP
highlights the dense spatial-description TR (``wooden'', ``chair'', ``propped''),
and LIME picks out the social-proximity words (\texttt{together},
\texttt{live}, \texttt{door}, \texttt{hand}) that dominate the narrative climax
window.

This cross-method agreement strengthens confidence that the identified content
reflects genuine neural selectivity rather than an artefact of either method's
design.

\subsubsection{Complementary Strengths}

The two methods provide distinct, non-redundant information:

\begin{itemize}
  \item \textbf{SHAP reveals latent BERT structure.}  By attributing importance
        to individual dimensions of the 768-dimensional BERT space, SHAP shows
        that brain-relevant semantic information is distributed across many
        dimensions (no single axis dominates) and is concentrated in early
        hemodynamic delays (delay block 1, 2\,s lag).  These findings hold
        globally across all stories and voxels and cannot be obtained from LIME,
        which operates on raw words rather than BERT features.

  \item \textbf{LIME reveals narrative-specific word salience.}  By linking
        response peaks to specific word tokens, LIME grounds the model's
        predictions in interpretable story content (mission dates, social
        actions) and enables comparison across subjects and adjacent voxels.
        SHAP word scores --- which assign the same value to all words in the
        same TR --- cannot distinguish between, say, \texttt{july} and
        \texttt{fourth} even though LIME shows them to have different
        contributions.

  \item \textbf{Response traces contextualise LIME windows.}  The predicted
        vs.\ observed traces (Figures~\ref{fig:resp_pluto}
        and~\ref{fig:resp_marry}) show that LIME's peak window corresponds to a
        moment of genuinely high predictive accuracy, validating that the local
        surrogate is explaining real signal rather than noise.  This cross-
        check is only possible because the encoding model provides continuous
        predictions that can be plotted against ground-truth BOLD.
\end{itemize}



\end{document}
